\documentclass[11pt,a4paper]{article}
\input{../../shared/preamble}
\SpeedHeader{Geometry}{3}

\begin{document}

\SpeedTitleBlock{Daily Geometry Drill \#3 --- Answers \& Investigations}{Henry Koduthore}
\noindent\textit{New toolkit today: Relative Positions $\cdot$ Radical Axis by subtraction $\cdot$ Common Chords $\cdot$ Orthogonal Circles $\cdot$ Tangency to Both Axes.}

%═══════════════════════════════════════════════════════════════════════════════
\section*{Section A \quad Rapid Recognition}
%═══════════════════════════════════════════════════════════════════════════════
\begin{enumerate}

%── A1 ──────────────────────────────────────────────────────────────────────────
\item The circle $\;x^2+y^2-12x-16y+75=0$ is tangent to the circle $\;x^2+y^2=25$. Write down the number of common tangents of the two circles.

\ans{$3$}
\method{$x^2+y^2-12x-16y+75=(x-6)^2+(y-8)^2-100+75=0$, so centre $(6,8)$, radius $\sqrt{36+64-75}=5$. Both circles have radius $5$; $d=10=r_1+r_2$ --- externally tangent --- $3$ common tangents.}
\inv{The bookkeeping: $d>r_1+r_2\Rightarrow4$, $d=r_1+r_2\Rightarrow3$, $|r_1-r_2|<d<r_1+r_2\Rightarrow2$, $d=|r_1-r_2|\Rightarrow1$, $d<|r_1-r_2|\Rightarrow0$. \textbf{``Tangent'' means exactly one common point, giving exactly $3$ tangents, not $2$.}}

%── A2 ──────────────────────────────────────────────────────────────────────────
\item The circles with centre $(0,0)$ radius $5$ and centre $(13,0)$ radius $3$ are disjoint. Write down the number of common tangents of the two circles.

\ans{$4$}
\method{$d=13$, $r_1+r_2=8$; since $d>r_1+r_2$ the circles are fully disjoint and there are $4$ common tangents.}
\inv{Disjoint circles admit $2$ direct (external) tangents plus $2$ transverse (internal) tangents. \textbf{Once the circles overlap, the two transverse tangents disappear.}}

%── A3 ──────────────────────────────────────────────────────────────────────────
\item Write down the radical axis of the circles $\;x^2+y^2-6x+2y+1=0$ and $\;x^2+y^2-2x-4y+13=0$.

\ans{"$2x-3y+6=0$"}
\method{At the radical axis powers are equal, so subtract: $(-6x+2y+1)-(-2x-4y+13)=0$, i.e.\ $-4x+6y-12=0$, i.e.\ $2x-3y+6=0$. The $x^2+y^2$ terms cancel in one stroke.}
\inv{Check the answer by symmetry: $x=\frac34 y-\frac32$ plug-in balances both powers automatically. \textbf{A single-line subtraction beats solving any intersection; the radical axis exists even for disjoint circles.}}

%── A4 ──────────────────────────────────────────────────────────────────────────
\item Write down whether the circles $\;x^2+y^2=9$ and $\;x^2+y^2+6x+8y+9=0$ are orthogonal (yes or no).

\ans{yes}
\method{$x^2+y^2+6x+8y+9$ has centre $(-3,-4)$ and $r^2=9+16-9=16$, so $r=4$. $d=\sqrt{3^2+4^2}=5$. Orthogonal iff $d^2=r_1^2+r_2^2$: $25=9+16$ \checkmark. Yes.}
\inv{Equivalently with coefficients: $g_1=f_1=0$, $g_2=3$, $f_2=4$ and $2g_1g_2+2f_1f_2=0$ must equal $c_1+c_2=(-9)+(9)=0$. \textbf{Orthogonality is about the radii forming a right triangle with the centre distance.}}

%── A5 ──────────────────────────────────────────────────────────────────────────
\item The circles $\;x^2+y^2=25$ and $\;x^2+y^2-10x+5=0$ intersect in two points. Find the length of their common chord.

\ans{$8$}
\method{Subtract: $-25=-10x+5\Rightarrow x=3$, the radical axis. Distance of $x=3$ from the origin is $3$, radius $5$, so half-chord $=\sqrt{25-9}=4$, chord $=8$.}
\inv{The second circle has centre $(5,0)$ and radius $\sqrt{20}$; the two circles agree on the two end-points of this chord. \textbf{Identify the radical axis first --- the chord sits on it --- then apply the half-chord formula.}}

%── A6 ──────────────────────────────────────────────────────────────────────────
\item Two circles have radii $3$ and $5$ and their centres are $4$ units apart. Write down the number of common tangents.

\ans{$2$}
\method{$|r_1-r_2|=2<d=4<r_1+r_2=8$ --- intersecting circles --- $2$ common tangents.}
\inv{Two crossing circles always have exactly $2$ common tangents. \textbf{The trap is counting the two intersection points as ``tangents''.}}

%── A7 ──────────────────────────────────────────────────────────────────────────
\item The circle $\;x^2+y^2-10x-10y+k=0$ is tangent to both coordinate axes. Find $k$.

\ans{$25$}
\method{Centre $(5,5)$. Tangent to both axes forces the radius to equal each coordinate: $r=5$. Then $r^2=25=5^2+5^2-k$, so $k=50-25=25$.}
\inv{A circle tangent to both axes has centre $(r,r)$, $(-r,r)$, $(r,-r)$ or $(-r,-r)$ and equation $(x\mp r)^2+(y\mp r)^2=r^2$, constant $r^2$. \textbf{The constant is $r^2$, not $2r^2$.}}

%── A8 ──────────────────────────────────────────────────────────────────────────
\item The circle $\;x^2+y^2+4x-2y+C=0$ is orthogonal to the circle $\;x^2+y^2-4x+4y-1=0$. Find $C$.

\ans{$-11$}
\method{$g_1=2$, $f_1=-1$; $g_2=-2$, $f_2=2$, $c_2=-1$. Orthogonality: $2g_1g_2+2f_1f_2=c_1+c_2\Rightarrow 2(2)(-2)+2(-1)(2)=C-1\Rightarrow-8-4=C-1\Rightarrow C=-11$.}
\inv{Geometric check: first circle centre $(-2,1)$ with $r_1^2=C$? No --- $r_1^2=4+1-C=5-C$; second centre $(2,-2)$ with $r_2^2=9$. $d^2=25=9+5-C\Rightarrow C=-11$ \checkmark. \textbf{Both forms of the test mirror each other.}}

%── A9 ──────────────────────────────────────────────────────────────────────────
\item How many common tangents do the circles $\;x^2+y^2-4x-6y+12=0$ and $\;x^2+y^2+2x-10y+25=0$ have?

\ans{$4$}
\method{First: centre $(2,3)$, $r^2=4+9-12=1$. Second: centre $(-1,5)$, $r^2=1+25-25=1$. $d=\sqrt{(3)^2+(2)^2}=\sqrt{13}>2=r_1+r_2$ --- disjoint unit circles --- $4$ common tangents.}
\inv{Completing the square (Worksheet \#2 skill) locates the centres; only then compare $d$ with $r_1+r_2$. \textbf{Worksheets \#2's centre-and-radius reading is today's first move.}}

%── A10 ─────────────────────────────────────────────────────────────────────────
\item Two circles have radii $7$ and $3$ and exactly $3$ common tangents. Find the distance between their centres.

\ans{$10$}
\method{Exactly $3$ common tangents means the circles are externally tangent, so $d=r_1+r_2=7+3=10$.}
\inv{Keep the five-case table on one line. \textbf{If there were $2$ tangents you would instead need $|3-7|<d<10$, which leaves a whole range, not a single value.}}

\end{enumerate}

%═══════════════════════════════════════════════════════════════════════════════
\section*{Section B \quad Manipulation Drills}
%═══════════════════════════════════════════════════════════════════════════════
\begin{enumerate}

%── B1 ──────────────────────────────────────────────────────────────────────────
\item How many common tangents do the circles $\;x^2+y^2=9$ and $\;x^2+y^2-4x-6y+9=0$ have?
\begin{itemize}
\item[A)] $0$
\item[B)] $1$
\item[C)] $2$
\item[D)] $3$
\item[E)] $4$
\end{itemize}

\ans{C) $2$}
\method{$x^2+y^2-4x-6y+9=0$: centre $(2,3)$, $r^2=4+9-9=4$. The other has centre $(0,0)$ radius $3$. $d=\sqrt{13}$; $|3-2|=1<\sqrt{13}<5$ --- intersecting --- $2$ common tangents.}
\inv{$x^2+y^2=9$ is radius $3$, so the second circle is strictly inside the first's $5$-radius relation. \textbf{Do not stop after finding the centres; compare $d$ against both $r_1+r_2$ and $|r_1-r_2|$.}}

%── B2 ──────────────────────────────────────────────────────────────────────────
\item The radical axis of the circles $\;x^2+y^2+ax+by+c=0$ and $\;x^2+y^2+6x-4y+10=0$ is $\;2x-y+3=0$. What is $a+b$?
\begin{itemize}
\item[A)] $3$
\item[B)] $-3$
\item[C)] $13$
\item[D)] $-13$
\end{itemize}

\ans{A) $3$}
\method{Subtract: $(a-6)x+(b+4)y+(c-10)=0$. Compare with $2x-y+3=0$: $a-6=2\Rightarrow a=8$; $b+4=-1\Rightarrow b=-5$; $c-10=3\Rightarrow c=13$. So $a+b=3$.}
\inv{Each coefficient matches with the {\em same} scale factor (here $1$); change the factor and all three coefficient equations change together. \textbf{The $y$-coefficient sign is the usual slip --- $b+4=-1$ gives $b=-5$, not $-3$.}}

%── B3 ──────────────────────────────────────────────────────────────────────────
\item The circle $\;x^2+y^2-4x+2y-5=0$ is orthogonal to $\;x^2+y^2+6x-2y+k=0$. Find $k$.
\begin{itemize}
\item[A)] $-9$
\item[B)] $9$
\item[C)] $-5$
\item[D)] $14$
\end{itemize}

\ans{A) $-9$}
\method{$g_1=-2,f_1=1,c_1=-5$; $g_2=3,f_2=-1,c_2=k$. $2g_1g_2+2f_1f_2=2(-2)(3)+2(1)(-1)=-14=c_1+c_2=-5+k\Rightarrow k=-9$.}
\inv{Orthogonal circles cut at right angles: the radii to a common point are perpendicular. \textbf{Never pair $g$ with $c$ or $f$ with $g$ --- the pairing is $g_1g_2+f_1f_2$ against $\frac{c_1+c_2}{2}$.}}

%── B4 ──────────────────────────────────────────────────────────────────────────
\item A circle is tangent to both coordinate axes in the first quadrant and passes through the point $(1,8)$. Find the smaller possible radius.

\ans{$5$}
\method{Centre $(r,r)$, radius $r$: $(1-r)^2+(8-r)^2=r^2\Rightarrow r^2-18r+65=0\Rightarrow(r-5)(r-13)=0$. Smaller radius $5$.}
\inv{There really are two such circles --- always state which root the question wants. \textbf{Both roots are positive, so the phrasing ``smaller'' is essential.}}

%── B5 ──────────────────────────────────────────────────────────────────────────
\item Two circles of equal radius $5$ have centres $(0,0)$ and $(6,0)$. Their radical axis has equation?
\begin{itemize}
\item[A)] $x=3$
\item[B)] $x=6$
\item[C)] $y=3$
\item[D)] $y=0$
\end{itemize}

\ans{A) $x=3$}
\method{Equal radii: subtract $(x^2+y^2-25)=((x-6)^2+y^2-25)\Rightarrow -25=-12x+36-25\Rightarrow x=3$.
}
\inv{For equal circles the radical axis is the perpendicular bisector of the centre segment. \textbf{A quick sanity check: power at $(3,0)$ is $9-25=-16$ for both sides.}}

%── B6 ──────────────────────────────────────────────────────────────────────────
\item Two circles of radii $4$ and $6$ have centres $9$ units apart. How many common tangents do they have?
\begin{itemize}
\item[A)] $0$
\item[B)] $1$
\item[C)] $2$
\item[D)] $3$
\item[E)] $4$
\end{itemize}

\ans{C) $2$}
\method{$|r_1-r_2|=2<d=9<10=r_1+r_2$: strict overlap without containment --- $2$ common tangents.}
\inv{Compare $d$ to $10$ (sum) and $2$ (difference) before counting. \textbf{$9$ is dangerously close to $10$; only exactly $10$ would give $3$.}}

%── B7 ──────────────────────────────────────────────────────────────────────────
\item The circle $\;x^2+y^2=4$ is orthogonal to $\;x^2+y^2+8x-6y+C=0$. Find $C$.
\begin{itemize}
\item[A)] $4$
\item[B)] $-4$
\item[C)] $25$
\item[D)] $2$
\end{itemize}

\ans{A) $4$}
\method{Second circle: centre $(-4,3)$, $r_2^2=25-C$. $d^2=16+9=25$. Orthogonal: $d^2=r_1^2+r_2^2=4+(25-C)=29-C\Rightarrow C=4$.}
\inv{Check with $g_1=f_1=0$: the coefficient test gives $0=-4+C\Rightarrow C=4$ \checkmark. \textbf{Both routes agree --- use whichever you trust.}}

%── B8 ──────────────────────────────────────────────────────────────────────────
\item The circles $\;x^2+y^2=36$ and $\;(x-6)^2+y^2=36$ intersect. Find the length of their common chord.

\ans{$6\sqrt{3}$}
\method{Equal radii $6$ on centres $6$ apart: radical axis is the perpendicular bisector $x=3$. Half-chord $=\sqrt{36-9}=3\sqrt{3}$, chord $=6\sqrt{3}$.}
\inv{The two circles overlap in a lens whose extreme width is $6\sqrt3$. \textbf{Do not add the two radii; the chord lives on the radical axis, not at the intersection of centre-lines.}}

%── B9 ──────────────────────────────────────────────────────────────────────────
\item A circle is tangent to both coordinate axes in the first quadrant and passes through the point $(4,8)$. Its possible radii are?
\begin{itemize}
\item[A)] $4$ only
\item[B)] $20$ only
\item[C)] $4$ and $20$
\item[D)] $16$
\end{itemize}

\ans{C) $4$ and $20$}
\method{Centre $(r,r)$: $(4-r)^2+(8-r)^2=r^2\Rightarrow r^2-24r+80=0\Rightarrow(r-4)(r-20)=0$. Both $4$ and $20$.}
\inv{Verify: radius $4$ circle touches both axes near the origin; radius $20$ circle is far out and still reaches $(4,8)$. \textbf{A quadratic in $r$ can give two genuine circles --- list both.}}

%── B10 ─────────────────────────────────────────────────────────────────────────
\item The point $P=(1,t)$ lies on the radical axis $\;3x-4y+5=0$ of two circles. Find $t$.

\ans{$2$}
\method{Substitute: $3(1)-4t+5=0\Rightarrow8-4t=0\Rightarrow t=2$.}
\inv{The radical axis is exactly the set of points of equal power --- a line. \textbf{Put $P$ into the line, not into either circle equation.}}

\end{enumerate}

%═══════════════════════════════════════════════════════════════════════════════
%═══════════════════════════════════════════════════════════════════════════════
\section*{Section C \quad Substitution \& Structure}
%═══════════════════════════════════════════════════════════════════════════════
\begin{enumerate}

%── C1 ──────────────────────────────────────────────────────────────────────────
\item \textit{\small(after TMUA 2021 Paper 1 Q1)} Two circles have the same radius. Their centres are $C_1(-2,1)$ and $C_2(3,-2)$. The circles intersect in two distinct points $P$ and $Q$. Which is the equation of the line $PQ$?
\begin{itemize}
\item[A)] $5x+3y=4$
\item[B)] $3x-5y=4$
\item[C)] $5x-3y=4$
\item[D)] $5x-3y=1$
\end{itemize}

\ans{C) $5x-3y=4$}
\method{Equal radii: subtract $(x+2)^2+(y-1)^2=(x-3)^2+(y+2)^2$, i.e.\ $x^2+4x+4+y^2-2y+1=x^2-6x+9+y^2+4y+4$, so $10x-6y=8$, i.e.\ $5x-3y=4$. The $x^2+y^2$ cancel in one line. The midpoint $(\tfrac12,-\tfrac12)$ satisfies $5(\tfrac12)-3(-\tfrac12)=4$ \checkmark.}
\inv{Equal radii make the radical axis the perpendicular bisector of $C_1C_2$, but you still get it fastest by subtraction. \textbf{Option A flips the sign of $y$; option D is a parallel line through the origin.}}

%── C2 ──────────────────────────────────────────────────────────────────────────
\item \textit{\small(after TMUA 2019 Paper 1 Q7)} The circles $C_1:(x-1)^2+y^2=4$ and $C_2:(x-5)^2+(y-3)^2=9$ have radii $2$ and $3$. How many common tangents do $C_1$ and $C_2$ have?
\begin{itemize}
\item[A)] $0$
\item[B)] $1$
\item[C)] $2$
\item[D)] $3$
\item[E)] $4$
\end{itemize}

\ans{D) $3$}
\method{$C_1=(1,0)$, $C_2=(5,3)$, $d=\sqrt{4^2+3^2}=5$. $r_1+r_2=5$, so $d=r_1+r_2$ --- externally tangent --- exactly $3$ common tangents ($2$ direct plus the common tangent at the touch point).}
\inv{$d>r_1+r_2\Rightarrow4$, $d=r_1+r_2\Rightarrow3$, $|r_1-r_2|<d<r_1+r_2\Rightarrow2$, $d=|r_1-r_2|\Rightarrow1$, $d<|r_1-r_2|\Rightarrow0$. \textbf{The $3$-$4$-$5$ triple makes $d$ instant.}}

%── C3 ──────────────────────────────────────────────────────────────────────────
\item \textit{\small(after TMUA 2016 Paper 2 Q9)} The circles $x^2+y^2=169$ and $x^2+y^2-6x+8y-119=0$ meet at $P$ and $Q$. What is the length $PQ$?
\begin{itemize}
\item[A)] $12$
\item[B)] $16$
\item[C)] $24$
\item[D)] $26$
\end{itemize}

\ans{C) $24$}
\method{Second circle: centre $(3,-4)$, $r^2=9+16+119=144$, $r=12$. Subtract: $-169=-6x+8y-119\Rightarrow6x-8y-50=0\Rightarrow3x-4y-25=0$. Distance from origin $=|{-}25|/5=5$. Half-chord $=\sqrt{169-25}=12$, chord $=24$. Check from $(3,-4)$: distance to line $|9-16-25|/5=1$, half-chord $\sqrt{144-1}= \sqrt{...}$ wait $r=12$, $d=1$ from second centre? No $d=|9+16-25|/5=0$? Let's recompute: $3(3)-4(-4)-25=9+16-25=0$? Actually $9+16-25=0$, so second centre lies on radical axis? Check again: line $3x-4y-25=0$ through $(3,-4)$ indeed $9+16-25=0$. So chord is diameter of second circle? No, check $r_2=12$, $r_1=13$, $d_{centres}=5$, not on axis? Wait calculation shows $3x-4y=25$ passes through $(3,-4)$? $9+16=25$ yes. That suggests second centre is on radical axis, which cannot be unless equal radii? Something off: original circle $x^2+y^2-6x+8y-119=0$ centre $(3,-4)$ radius $\sqrt{9+16+119}=12$, first circle origin radius13. Radical axis is $6x-8y=50$, i.e.\ $3x-4y=25$. Plug $(3,-4)$: $9+16=25$ true, so second centre lies on axis — that would mean power of second centre is same? But power of centre w.r.t own circle is $-r^2=-144$, w.r.t first circle is $9+16-169=-144$ same, so indeed second centre has equal power because $d^2=25$, $r_1^2-r_2^2=169-144=25$, so radical axis passes through second centre's projection? Actually condition $d^2=r_1^2-r_2^2$ makes axis pass through second centre, possible when $d^2 = r_1^2 - r_2^2$ gives one centre on axis. Here $25=25$ holds, so yes axis passes through $(3,-4)$. Then distance from second centre to axis is $0$, chord would be diameter $24$ of second circle, consistent! So half-chord from origin $12$ matches $r_2=12$ diameter. Both views agree: chord $=24$.}
\inv{The $5$-$12$-$13$ triple runs the whole question. Option D $26$ is the diameter of the large circle, strictly longer than any chord. \textbf{The radical axis can pass through a centre when $d^2=|r_1^2-r_2^2|$.}}

%── C4 ──────────────────────────────────────────────────────────────────────────
\item \textit{\small(adapted from JZMaths Mock 1 Q5)} A circle in the first quadrant is tangent to both coordinate axes and passes through the point $(9,2)$. Its two possible radii are $r_1<r_2$. What is $r_1+r_2$?
\begin{itemize}
\item[A)] $14$
\item[B)] $17$
\item[C)] $22$
\item[D)] $34$
\end{itemize}

\ans{C) $22$}
\method{Centre $(r,r)$, radius $r$: $(9-r)^2+(2-r)^2=r^2\Rightarrow r^2-22r+85=0$. By Vieta sum $=22$, product $=85$; roots are $5$ and $17$. Sum $22$.}
\inv{Both $r=5$ and $r=17$ give genuine circles. \textbf{Option B $17$ is the larger root alone; option A $14$ would be $9+5$ misreading.}}

%── C5 ──────────────────────────────────────────────────────────────────────────
\item \textit{\small(after Oxbridge Paper 1 Q8)} The circle $x^2+y^2+4x-2y+C=0$ is orthogonal to the circle $x^2+y^2-4x+4y-1=0$. What is $C$?
\begin{itemize}
\item[A)] $-11$
\item[B)] $-8$
\item[C)] $11$
\item[D)] $5$
\end{itemize}

\ans{A) $-11$}
\method{$g_1=2$, $f_1=-1$; $g_2=-2$, $f_2=2$, $c_2=-1$. Orthogonality $2g_1g_2+2f_1f_2=c_1+c_2\Rightarrow2(2)(-2)+2(-1)(2)=C-1\Rightarrow-8-4=C-1\Rightarrow C=-11$. Geometric check: centres $(-2,1)$ and $(2,-2)$, $d^2=25$, $r_1^2=5-C$, $r_2^2=9$, $25=5-C+9\Rightarrow C=-11$ \checkmark.}
\inv{Pair $g$ with $g$, $f$ with $f$: $2g_1g_2+2f_1f_2$ against $(c_1+c_2)$. \textbf{Option C $11$ flips the sign of the constant.}}

%── C6 ──────────────────────────────────────────────────────────────────────────
\item \textit{\small(after TMUA 2021 Paper 2 Q15)} A circle has equation $x^2+ax+y^2+by+c=0$ where $a,b,c\neq0$. Which condition is necessary and sufficient for the circle to be tangent to the $y$-axis?
\begin{itemize}
\item[A)] $a^2=4c$
\item[B)] $b^2=4c$
\item[C)] $a^2+b^2=4c$
\item[D)] $c=0$
\end{itemize}

\ans{B) $b^2=4c$}
\method{Centre $(-a/2,-b/2)$, $r^2=a^2/4+b^2/4-c$. Tangent to $y$-axis ($x=0$) needs distance $|{-}a/2|=r$, so $a^2/4=a^2/4+b^2/4-c\Rightarrow b^2=4c$. (Tangency to $x$-axis would be $a^2=4c$.)}
\inv{The $y$-axis condition kills $b$, not $a$, because the radius condition cancels $a^2/4$. \textbf{Option A $a^2=4c$ is tangency to the $x$-axis; option D $c=0$ is through the origin.}}

%── C7 ──────────────────────────────────────────────────────────────────────────
\item \textit{\small(adapted from TMUA 2020 Paper 1 Q16)} The equal circles $C_1:(x+2)^2+(y-1)^2=3$ and $C_2:(x-4)^2+(y-1)^2=3$ have centres $6$ apart. A transverse common tangent with positive gradient makes an acute angle $\theta$ with the $x$-axis. What is $\tan\theta$?
\begin{itemize}
\item[A)] $\dfrac12$
\item[B)] $\dfrac{\sqrt2}{2}$
\item[C)] $\sqrt2$
\item[D)] $\dfrac{\sqrt3}{3}$
\end{itemize}

\ans{B) $\dfrac{\sqrt2}{2}$}
\method{Equal radii $\sqrt3$ at $(-2,1)$ and $(4,1)$. Transverse tangent passes through midpoint $(1,1)$: line $y-1=m(x-1)$, i.e.\ $mx-y+1-m=0$. Distance from $(-2,1)$ to line $=|{-}3m|/\sqrt{m^2+1}=\sqrt3\Rightarrow9m^2=3(m^2+1)\Rightarrow6m^2=3\Rightarrow m^2=1/2\Rightarrow m=\sqrt2/2$ (positive).}
\inv{The transverse tangent crosses the centre line at the midpoint; the direct common tangents are horizontal ($m=0$). \textbf{Option C $\sqrt2$ inverts the fraction; option A $1/2$ squares the root.}}

%── C8 ──────────────────────────────────────────────────────────────────────────
\item \textit{\small(adapted from Vantage Mock Paper 2 Q10 $/$ BMO1 2017 Q4)} The three circles $C_1:x^2+y^2+4x-1=0$, $C_2:x^2+y^2-8y+11=0$, $C_3:x^2+y^2+2x-12y+13=0$ have pairwise radical axes concurrent at $R$. What are the coordinates of $R$?
\begin{itemize}
\item[A)] $(1,1)$
\item[B)] $(1,-1)$
\item[C)] $(3,1)$
\item[D)] $(1,2)$
\end{itemize}

\ans{A) $(1,1)$}
\method{Radical axis $C_1-C_2$: $(4x-1)-(-8y+11)=0\Rightarrow4x+8y-12=0\Rightarrow x+2y-3=0$. $C_2-C_3$: $(-8y+11)-(2x-12y+13)=0\Rightarrow-2x+4y-2=0\Rightarrow x-2y+1=0$. Solve $x+2y=3$, $x-2y=-1\Rightarrow x=1,y=1$. Power: $1+1+4-1=5$ for all three \checkmark.}
\inv{Any two radical axes intersect at the radical centre; the third must pass through it. \textbf{Option D $(1,2)$ lies on only one axis.}}

\end{enumerate}

%═══════════════════════════════════════════════════════════════════════════════
\section*{Section D \quad Challenge Ramp}
%═══════════════════════════════════════════════════════════════════════════════
\begin{enumerate}

%── D1 ──────────────────────────────────────────────────────────────────────────
\item \textit{\small(adapted from TMUA 2018 Paper 1 Q3)} The circles $C_1:(x+2)^2+(y-3)^2=18$ and $C_2:(x-7)^2+(y+6)^2=2$ have centres $O_1$ and $O_2$. What is the shortest distance between a point on $C_1$ and a point on $C_2$?
\begin{itemize}
\item[A)] $5\sqrt2-4$
\item[B)] $5\sqrt2-5$
\item[C)] $5\sqrt2$
\item[D)] $9\sqrt2$
\end{itemize}

\ans{C) $5\sqrt2$}
\method{$O_1=(-2,3)$, $O_2=(7,-6)$, $O_1O_2=\sqrt{9^2+9^2}=9\sqrt2$, $r_1=\sqrt{18}=3\sqrt2$, $r_2=\sqrt2$, sum $4\sqrt2$, shortest $=9\sqrt2-4\sqrt2=5\sqrt2$. Externally disjoint ($d>r_1+r_2$).}
\inv{Shortest $=d-r_1-r_2$ when disjoint; if one contained the other it would be $|r_1-r_2|-d$. \textbf{Option D $9\sqrt2$ is $O_1O_2$ itself, forgetting to subtract radii.}}

%── D2 ──────────────────────────────────────────────────────────────────────────
\item \textit{\small(after TMUA 2019 Paper 1 Q6)} The circles $(x+4)^2+(y+1)^2=64$ and $(x-8)^2+(y-4)^2=r^2$ ($r>0$) have exactly one point in common. What is the difference between the two possible values of $r$?
\begin{itemize}
\item[A)] $8$
\item[B)] $16$
\item[C)] $26$
\item[D)] $42$
\end{itemize}

\ans{B) $16$}
\method{Centre distance $d=\sqrt{12^2+5^2}=13$, $r_1=8$. One point $\Rightarrow|r-8|=13$ (internal) or $r+8=13$ (external) $\Rightarrow r=21$ or $5$, difference $16$.}
\inv{External tangency gives smaller $r=5$, internal gives larger $21$. \textbf{Option A $8$ is $r_1$ itself; option D $42$ would be sum not difference.}}

%── D3 ──────────────────────────────────────────────────────────────────────────
\item \textit{\small(adapted from BMO1 2017 Q4)} Three circles of radii $3$, $4$ and $5$ are pairwise externally tangent. What is the area of the triangle formed by joining their centres?
\begin{itemize}
\item[A)] $12\sqrt5$
\item[B)] $12$
\item[C)] $84$
\item[D)] $30$
\end{itemize}

\ans{A) $12\sqrt5$}
\method{Centre triangle sides $3+4=7$, $3+5=8$, $4+5=9$. $s=12$, Heron area $=\sqrt{12\cdot5\cdot4\cdot3}=\sqrt{720}=12\sqrt5$. Coordinate check: $(0,0)$, $(7,0)$, third $x=33/7$, $y=24\sqrt5/7$, area $=\tfrac12\cdot7\cdot24\sqrt5/7=12\sqrt5$ \checkmark.}
\inv{Side is sum of radii, not difference. \textbf{Option C $84$ is area of $13$-$14$-$15$ triangle, the classic Heron triple.}}

%── D4 ──────────────────────────────────────────────────────────────────────────
\item \textit{\small(after TMUA 2016 Paper 2 Q9 $/$ Oxbridge)} Two circles of radii $3$ and $4$ are orthogonal. What is the length of their common chord?
\begin{itemize}
\item[A)] $\dfrac{24}{5}$
\item[B)] $\dfrac{12}{5}$
\item[C)] $6$
\item[D)] $\dfrac{48}{5}$
\end{itemize}

\ans{A) $\dfrac{24}{5}$}
\method{$d^2=r_1^2+r_2^2=9+16=25$, $d=5$. Place centres $(0,0)$ and $(5,0)$; radical axis $10x=18\Rightarrow x=9/5$ from small centre. Half-chord $\sqrt{9-(9/5)^2}=\sqrt{144/25}=12/5$, chord $=24/5$. Also $chord=2r_1r_2/d=24/5$.}
\inv{Option C $6$ is diameter of small circle; chord is strictly shorter. \textbf{Option B $12/5$ is half the chord.}}

%── D5 ──────────────────────────────────────────────────────────────────────────
\item \textit{\small(adapted from JZMaths $/$ Vantage)} A point $P$ has equal power with respect to the circles $x^2+y^2=1$ and $(x-6)^2+y^2=4$. What does the locus of $P$ describe?
\begin{itemize}
\item[A)] the line $x=\dfrac{11}{4}$
\item[B)] the circle $x^2+y^2=1$
\item[C)] the $y$-axis
\item[D)] a single point $\left(\dfrac{11}{4},0\right)$
\end{itemize}

\ans{A) the line $x=\dfrac{11}{4}$}
\method{Equal powers $=$ radical axis: $x^2+y^2-1=(x-6)^2+y^2-4=x^2-12x+32+y^2\Rightarrow12x=33\Rightarrow x=11/4$, a vertical line. Exists even though circles are disjoint ($d=6>3$).}
\inv{Locus of equal power is always a straight line, even for disjoint circles. \textbf{Option D is a single point on that line, not the whole locus.}}

\end{enumerate}

%═══════════════════════════════════════════════════════════════════════════════
\section*{Top 5 Patterns --- Worksheet \#3}
\begin{enumerate}[leftmargin=2em, itemsep=4pt]
  \item \textbf{Classify two circles by comparing $d$ with $r_1+r_2$ and $|r_1-r_2|$.} $d>r_1+r_2$ (4 tangents), $d=r_1+r_2$ (3), overlap (2), $d=|r_1-r_2|$ (1), containment (0). \seealso{A1, A2, A6, B1, B6, C2, D1, D2}
  \item \textbf{Radical axis $=$ the two equations subtracted.} The $x^2+y^2$ terms vanish in one line; the result is the common chord for intersecting circles and the equal-power line in general. \seealso{A3, B2, C1, C3, C8, D5}
  \item \textbf{Common chord $=$ radical axis $+$ half-chord formula.} Find the radical axis, take its distance $d$ from a centre, then $2\sqrt{r^2-d^2}$. \seealso{A5, B8, C3, D4}
  \item \textbf{Orthogonality either way: $d^2=r_1^2+r_2^2$ or $2g_1g_2+2f_1f_2=c_1+c_2$.} The radii make a right triangle against the centre distance. \seealso{A4, A8, B3, B7, C5, D4}
  \item \textbf{A circle tangent to both axes lives at $(\pm r,\pm r)$ with constant $r^2$.} Powers of $r$ are found from a point condition; Vieta gives sum and product instantly. \seealso{A7, B4, B9, C4}
\end{enumerate}

\section*{Common Traps --- Worksheet \#3}
\begin{enumerate}[leftmargin=2em, itemsep=4pt]
  \item \textbf{Miscounting the common-tangent table.} Externally tangent means exactly $3$ common tangents (diameter line plus two direct); crossing circles have $2$, disjoint circles $4$. One boundary-off-by-one changes the whole count. \seealso{A1, A2, A6, A10, B1, B6, C2, D1, D2}
  \item \textbf{Mangling the subtraction sign on the radical axis.} It is the difference of the two {\em constants and linear terms} set to zero; flip a sign and you get the next option in C1 or C5. \seealso{A3, B2, C1, C5, C8}
  \item \textbf{Wrong pairing in the orthogonality test.} The condition binds $g_1g_2+f_1f_2$ against $(c_1+c_2)/2$ --- pair $g$ with $g$, $f$ with $f$. \seealso{A8, B3, B7, C5, D4}
  \item \textbf{Imagining the radical axis is a circle.} Equal powers trace a straight line, not a curve, even between disjoint circles. \seealso{B5, B10, C1, D5}
  \item \textbf{Forgetting the triangle of centres on tangent circles.} For pairwise tangent circles each centre distance is the {\em sum} of the two radii; subtracting instead of adding breaks Heron immediately. \seealso{D3, C2, A10}
\end{enumerate}

\SpeedClosing{``Two circles are solved by one subtraction: locate, subtract, then compare with the sum and difference of the radii.''}

\end{document}
